\section*{Capítulo \ref{ch:b2:speciation} --- Especiación y macroevolución}

\begin{solution}{exo:b2:speciation:1}
Una especie es un grupo de poblaciones naturales que se cruzan entre sí y
que están aisladas reproductivamente de los demás grupos de ese tipo.
Falla para los linajes asexuales (bacterias, rotíferos bdeloideos, dientes
de león: concepto morfológico o filogenético, conjuntos de formas
parecidas o linajes diagnosticables); para los fósiles (concepto
morfológico, puesto que no puede hacerse ningún cruzamiento); y para las
poblaciones alopátricas que nunca se encuentran (concepto ecológico o
filogenético, puesto que el cruzamiento no puede ponerse a prueba) --- y lo
difuminan los robles que se hibridan y las especies en anillo.
\end{solution}

\begin{solution}{exo:b2:speciation:2}
Ranas en meses distintos: precigótico (temporal). Mula: poscigótico
(esterilidad híbrida). Espermatozoide que no puede unirse: precigótico
(gamético). Híbridos vigorosos con \omterm{def:b2:angiosperm-reproduction:seed}{semillas} estériles: poscigótico
(esterilidad híbrida, o degradación híbrida si el fallo se produce en la
generación siguiente). Grillos con cantos distintos: precigótico
(comportamental).
\end{solution}

\begin{solution}{exo:b2:speciation:3}
Alopátrica: una barrera divide un área de distribución (los camarones
pistola de los dos lados del istmo de Panamá). Peripátrica: unos pocos
fundadores más allá del área (los pinzones de Darwin a partir de un
antepasado continental). Simpátrica: sin barrera (la mosca de la manzana
sobre el espino albar y el manzano; un trigo \omterm{def:b2:genome-diversification:polyploidy}{alopoliploide}). Parapátrica:
a lo largo de un gradiente, con una \omterm{prop:b2:speciation:reinforcement}{zona híbrida} en el centro (gramíneas
dentro y fuera de escombreras de mina, que florecen en épocas distintas).
\end{solution}

\begin{solution}{exo:b2:speciation:4}
Gana siempre la especie 1; gana siempre la especie 2; coexistencia
estable; gana una u otra según los efectivos iniciales. La coexistencia es
estable cuando cada especie puede invadir a la otra en su capacidad de
carga, es decir, cuando cada especie se limita a sí misma más de lo que
limita a la otra ($K_1 > \alpha K_2$ y $K_2 > \beta K_1$).
\end{solution}

\begin{solution}{exo:b2:speciation:5}
$\alpha K_2 = 400 < K_1 = 1000$: la especie 1 invade; $\beta K_1 = 600
< K_2 = 800$: la especie 2 invade. $N_1^{*} = (1000 - 400)/(1 - 0.3) =
857$, $N_2^{*} = (800 - 600)/0.7 = 286$; en el equilibrio, la especie
1 está por debajo de su $K$ en $\alpha N_2^{*} = 143$ y la especie 2 en $\beta
N_1^{*} = 514$.
\end{solution}

\begin{solution}{exo:b2:speciation:6}
$\alpha K_2 = 1200 > K_1$: la especie 1 no puede invadir a la especie 2 en
su capacidad de carga; $\beta K_1 = 300 < K_2$: la especie 2 puede invadir
a la especie 1. Gana siempre la especie 2, porque sus individuos
perjudican a la especie 1 más que los propios de la especie 1, mientras
que ella apenas se ve afectada. Con $\alpha = \beta = 1.5$: $\alpha K_2 = 1200 > 1000$
y $\beta K_1 = 1500 > 800$: ninguna puede invadir a la otra, cada una
excluye a la otra una vez establecida, y gana la que llegó primero.
\end{solution}

\begin{solution}{exo:b2:speciation:7}
El tetraploide $4n$ produce gametos $2n$; cruzado con los gametos $n$ del
\omterm{def:b2:meiosis-heredity:meiosis}{diploide} da triploides $3n$, cuya \omterm{def:b2:meiosis-heredity:meiosis}{meiosis} no puede emparejar tres juegos y
produce gametos aneuploides, en su mayoría inviables. El tetraploide solo
puede reproducirse con otros tetraploides (o por autofecundación): queda
aislado de un solo paso. Suele fracasar porque está solo --- sus únicas
parejas posibles son \omterm{def:b2:meiosis-heredity:meiosis}{diploides} y cada uno de esos cruzamientos se
desperdicia --- y queda sumergido; persiste cuando se autofecunda, se
propaga vegetativamente o surge repetidas veces.
\end{solution}

\begin{solution}{exo:b2:speciation:8}
$0.01\times 10\times 10 = 1$ incompatibilidad; con 30 en cada una, $0.01\times
900 = 9$. El aislamiento crece como el cuadrado del tiempo de divergencia:
dos poblaciones que han divergido durante un tiempo tres veces mayor están
nueve veces más aisladas --- la bola de nieve.
\end{solution}

\begin{solution}{exo:b2:speciation:9}
$w = 2\sqrt{8/0.05} = \qty{25}{km}$; $w = 0.3\sqrt{8/0.4} =
\qty{1.3}{km}$. El refuerzo es más probable en la segunda: los híbridos
son muy poco aptos, de modo que una hembra que rechaza a la otra forma
gana mucho, y la zona es lo bastante estrecha para que las dos formas se
encuentren con la frecuencia suficiente para que ese rechazo se
seleccione.
\end{solution}

\begin{solution}{exo:b2:speciation:10}
$\alpha K_2 = 450 < 500$, de modo que cada una invade a la otra: coexisten,
en $N^{*} = (500 - 450)/(1 - 0.81) = 263$ cada una --- por poco, cerca del
límite de la exclusión, y un pequeño cambio de $K$ inclinaría la balanza.
Con $\alpha = \beta = 0.4$: $N^{*} = (500 - 200)/(1 - 0.16) = 357$ cada
una. El desplazamiento ha girado cada isoclina alejándola de la otra
($K/\alpha$ se ha desplazado hacia fuera), de modo que el cruce se ha
alejado de los ejes y cada especie se sitúa más cerca de su propia $K$:
una coexistencia más segura y con más individuos de cada una.
\end{solution}

\begin{solution}{exo:b2:speciation:11}
$R = h^{2}S = 0.7\times(-1) = \qty{-0.7}{mm}$: lo observado. En la sequía
de 1977, sin el pinzón terrestre grande, los mismos pinzones medianos
fueron seleccionados a favor de picos \emph{más grandes}, puesto que las
\omterm{def:b2:angiosperm-reproduction:seed}{semillas} grandes eran el alimento que quedaba; en 2004, las \omterm{def:b2:angiosperm-reproduction:seed}{semillas}
grandes las acaparaba el invasor, y los pinzones medianos que competían
por ellas perdieron: el sentido de la selección lo fijó el competidor.
\end{solution}

\begin{solution}{exo:b2:speciation:12}
Se hacen por accidente: el modelo de dos loci muestra que el aislamiento
surge como subproducto de sustituciones fijadas por otras razones, o por
deriva, y la geografía --- una barrera, una fundación --- proporciona la
separación que permite que el accidente se acumule. Se mantienen por
selección: donde las formas se encuentran, la selección contra los
híbridos poco aptos construye barreras precigóticas (refuerzo), y la
competencia desplaza los caracteres para que ambas puedan coexistir. El
eslogan omite la \omterm{prop:b2:speciation:modes}{especiación simpátrica} por selección disruptiva, en la
que la selección hace además de mantener las especies, y la \omterm{def:b2:genome-diversification:polyploidy}{poliploidía},
en la que el accidente lo es todo.
\end{solution}

\begin{solution}{pb:b2:speciation:1}
\textbf{1.} Especie 1: $N_1 + 0.9N_2 = 1200$, cortes $N_1 = 1200$
y $N_2 = 1333$. Especie 2: $N_2 + 0.3N_1 = 400$, cortes $N_2 =
400$ y $N_1 = 1333$.
\textbf{2.} $\beta K_1 = 360 < K_2 = 400$: la especie 2 invade;
$\alpha K_2 = 360 < K_1 = 1200$: la especie 1 invade. Coexistencia
estable.
\textbf{3.} $N_1^{*} = (1200 - 360)/(1 - 0.27) = 1151$; $N_2^{*} =
(400 - 360)/0.73 = 55$.
\textbf{4.} $\mathrm{d}N_1/\mathrm{d}t = 0.5\times 1200\,(1 - 1218/1200)
= -9$ por año: el residente, que estaba en su capacidad de carga, empieza a
disminuir empujado por los recién llegados.
\textbf{5.} $K_1 = 600$, $K_2 = 200$: $N_1^{*} = (600 - 180)/0.73 =
575$, $N_2^{*} = (200 - 180)/0.73 = 27$: ambos se reducen a la mitad, y el
invasor, con 27 aves, está cerca de la extinción --- la sequía es cuando
la competencia aprieta.
\textbf{6.} El invasor es el ave más grande, de pico más grande: toma las
\omterm{def:b2:angiosperm-reproduction:seed}{semillas} grandes que el residente también usa, y las toma más deprisa,
mientras que las \omterm{def:b2:angiosperm-reproduction:seed}{semillas} pequeñas del residente le sirven de poco.
\textbf{7.} $S = 9.0 - 9.5 = \qty{-0.5}{mm}$; $R = 0.7\times(-0.5) =
\qty{-0.35}{mm}$.
\textbf{8.} $9.5 - 0.35 = \qty{9.15}{mm}$.
\textbf{9.} $N_1^{*} = (1200 - 200)/(1 - 0.15) = 1176$; $N_2^{*} =
(400 - 360)/0.85 = 47$.
\textbf{10.} El corte con el eje $N_2$ de la isoclina de la especie 1 ha
subido de
$K_1/0.9 = 1333$ a $K_1/0.5 = 2400$: la recta ha girado alejándose de la
del invasor, el residente se ve ahora menos afectado por cada invasor, y el
equilibrio se ha desplazado hacia la $K$ del residente.
\textbf{11.} $2/0.35 \approx 6$ años. No ocurre así porque una selección de
este tipo solo se da en los años de sequía, se invierte en los años húmedos,
cuando abundan las \omterm{def:b2:angiosperm-reproduction:seed}{semillas} pequeñas (o cuando el invasor escasea), y
porque la varianza aditiva se agotaría.
\textbf{12.} Dos millones de años son alrededor de $10^{6}$ generaciones;
un desplazamiento de \qty{0.35}{mm} por generación sostenido durante tan
solo una milésima de ellas cambiaría el pico en cientos de milímetros. La
selección es fuerte pero inconstante; el cambio direccional sostenido es
la rara excepción, y el cambio neto a escala geológica es un residuo
diminuto de grandes episodios opuestos.
\textbf{13.} El canto, aprendido del padre y usado por las hembras en la
elección de pareja, con la propia forma del pico como señal. Los híbridos
de pico intermedio no manejan bien ni las \omterm{def:b2:angiosperm-reproduction:seed}{semillas} grandes ni las
pequeñas, y son los primeros en morir de hambre durante las sequías.
\textbf{14.} $w = 0.5\sqrt{8/0.2} = \qty{3.2}{km}$.
\textbf{15.} La isla es más estrecha de lo que sería la zona: no puede
formarse ninguna clina espacial; la hibridación, donde se produce, afecta
a toda la población, y solo el apareamiento selectivo mantiene separadas
a las dos especies.
\textbf{16.} Las hembras que solo se aparean con machos de su propio canto
y su propio pico dejan más nietos que las que se hibridan; la preferencia
y la señal (canto, pico) se vuelven más distintas donde las dos especies
viven juntas que donde viven solas.
\textbf{17.} $0.005\times 15\times 15 = 1.1$ incompatibilidades: una por
término medio, con híbridos probablemente disminuidos pero no estériles.
Todavía no son especies según el criterio biológico; están en camino.
\textbf{18.} Biológico: todavía no (híbridos casi aptos). Morfológico:
quizá, si los picos o el plumaje han divergido de forma visible.
Filogenético: sí, si cada población tiene una diferencia diagnóstica
fijada. Los conceptos discrepan precisamente durante el proceso de
especiación, y eso es lo que los hace útiles.
\textbf{19.} $1/(5\times 10^{6}) = 2\times 10^{-7}$ por especie y por
año; para $10^{7}$ especies, dos extinciones por año.
\textbf{20.} De 200 a 2000 especies por año.
\textbf{21.} $0.75\times 10^{7}/10^{5} = 75$ especies por año --- menos
que la tasa actual, incluso en su valor bajo. El episodio actual, si se
mantiene, es una \omterm{prop:b2:speciation:tempo}{extinción masiva} que avanza más deprisa que las cinco
grandes.
\textbf{22.} Los nichos que dejan libres los perdedores están vacíos de
competidores, de modo que cualquier variante superviviente capaz de usar
uno queda favorecida, y los supervivientes se diversifican hacia ellos.
\textbf{23.} \omterm{prop:b2:speciation:tempo}{Equilibrio puntuado}; \omterm{prop:b2:speciation:modes}{especiación alopátrica} (peripátrica) en
una pequeña población aislada cuyo descendiente se extiende después por el
área de la población principal.
\textbf{24.} $1.04^{n} = 2$: $n = \ln 2/\ln 1.04 = 18$ generaciones.
El registro muestra millones de años porque la selección se invierte, la
\omterm{prop:b2:population-genetics:kinds}{selección estabilizadora} domina la mayor parte del tiempo y el cambio neto
es la pequeña diferencia entre grandes episodios opuestos.
\textbf{25.} Antes del desplazamiento: coexistencia estable con $N_1^{*} =
1151$, $N_2^{*} = 55$; después: 1176 y 47. Desplazamiento del pico de
\qty{-0.35}{mm} en una generación. Anchura de la \omterm{prop:b2:speciation:reinforcement}{zona híbrida}:
\qty{3.2}{km}.
\end{solution}
